\usepackage[margin=1in]{geometry}
\usepackage{float}
\interfootnotelinepenalty=10000
\usepackage{wrapfig}
\usepackage{titlesec}
\titlespacing*{\section}{0pt}{9pt}{4pt}
\titlespacing*{\subsection}{0pt}{6pt}{3pt}
% let pandoc figures float (htbp) so text fills page bottoms instead of leaving gaps
\let\origfigure\figure
\let\endorigfigure\endfigure
\renewenvironment{figure}[1][]{%
  \origfigure[htbp]
}{%
  \endorigfigure
}
% allow floats to occupy more of a page, reducing float-induced whitespace
\renewcommand{\topfraction}{0.9}
\renewcommand{\bottomfraction}{0.9}
\renewcommand{\textfraction}{0.07}
\renewcommand{\floatpagefraction}{0.8}
\usepackage{caption}
\captionsetup{font=small, labelfont=bf, labelformat=empty}
\setlength{\abovecaptionskip}{4pt}
\setlength{\belowcaptionskip}{0pt}
\usepackage{graphicx}
\setkeys{Gin}{width=0.85\linewidth,keepaspectratio}
% let the Markdown use Unicode math/relation symbols freely
% colored placeholder-chit boxes (see chit.lua): amber = figure, blue = table
\usepackage[most]{tcolorbox}
\newtcolorbox{placeholderfig}{breakable, enhanced, colback=orange!13, colframe=orange!80!black,
  boxrule=1.2pt, arc=2pt, left=6pt, right=6pt, top=5pt, bottom=5pt}
\newtcolorbox{placeholdertable}{breakable, enhanced, colback=blue!8, colframe=blue!65!black,
  boxrule=1.2pt, arc=2pt, left=6pt, right=6pt, top=5pt, bottom=5pt}
\usepackage{newunicodechar}
\newunicodechar{≥}{\ensuremath{\geq}}
\newunicodechar{≤}{\ensuremath{\leq}}
\newunicodechar{≈}{\ensuremath{\approx}}
\newunicodechar{≡}{\ensuremath{\equiv}}
\newunicodechar{→}{\ensuremath{\rightarrow}}
\newunicodechar{↔}{\ensuremath{\leftrightarrow}}
\newunicodechar{×}{\ensuremath{\times}}
\newunicodechar{∼}{\ensuremath{\sim}}
\newunicodechar{−}{\ensuremath{-}}
\newunicodechar{⁰}{\ensuremath{^{0}}}
\newunicodechar{¹}{\ensuremath{^{1}}}
\newunicodechar{²}{\ensuremath{^{2}}}
\newunicodechar{³}{\ensuremath{^{3}}}
\newunicodechar{⁴}{\ensuremath{^{4}}}
\newunicodechar{⁵}{\ensuremath{^{5}}}
\newunicodechar{⁶}{\ensuremath{^{6}}}
\newunicodechar{⁷}{\ensuremath{^{7}}}
\newunicodechar{⁸}{\ensuremath{^{8}}}
\newunicodechar{⁹}{\ensuremath{^{9}}}
\newunicodechar{σ}{\ensuremath{\sigma}}
\newunicodechar{α}{\ensuremath{\alpha}}
